\chapter{Engenharia de Requisitos}
\label{cap_requisitos}

A Engenharia de Requisitos é a disciplina central da Engenharia de Software responsável por descobrir, refinar, modelar e documentar os serviços exigidos e as restrições operacionais de um sistema \cite{sommerville2019}. No SIGEA-GTT, os requisitos foram elicitados a partir das normas do IHI \cite{griffin2009ihi}, das diretrizes do PNSP/ANVISA \cite{anvisa2017} e dos consensos didático-pedagógicos estabelecidos entre o Departamento de Computação (DCOMP) e o Departamento de Enfermagem da Universidade Federal de Sergipe (UFS).

\section{Atores do Sistema}
\label{sec_atores}

O modelo de controle de acesso do SIGEA-GTT é estritamente baseado em perfis funcionais (Role-Based Access Control - RBAC):
\begin{enumerate}
    \item \textbf{Administrador (\texttt{ADMIN})}: Responsável pela administração do ecossistema, incluindo cadastro de usuários, parametrização do catálogo de módulos e gatilhos clínicos do IHI, manutenção das categorias de gravidade NCC MERP e criação e supervisão de turmas acadêmicas;
    \item \textbf{Professor (\texttt{PROFESSOR})}: Docente responsável por gerenciar suas turmas acadêmicas, elaborar atividades práticas com prontuários simulados de alta fidelidade clínica, acompanhar as submissões dos estudantes, realizar a avaliação pedagógica com atribuição de nota e parecer formativo, e analisar os dashboards de indicadores da turma;
    \item \textbf{Estudante (\texttt{STUDENT})}: Discente de graduação em saúde (enfermagem, medicina e áreas afins) que consulta suas turmas matriculadas, realiza a auditoria retrospectiva dos prontuários clínicos simulados sob o cronômetro de 20 minutos, identifica gatilhos e danos, aplica as seis Ferramentas da Qualidade, submete sua resolução e consulta o parecer formativo emitido pelo docente.
\end{enumerate}

\section{Requisitos Funcionais (RF01 a RF25)}
\label{sec_requisitos_funcionais}

Os Requisitos Funcionais descrevem as funções, comportamentos e regras de negócio específicas que o software deve executar. As Tabelas \ref{tab_rf_1_13} e \ref{tab_rf_14_25} consolidam os 25 requisitos funcionais implementados no SIGEA-GTT, identificando seus códigos, descrições detalhadas, prioridades (Essencial ou Importante) e atores autorizados.

\begin{table}[htb]
\centering
\small
\caption{Especificação dos Requisitos Funcionais (RF01 a RF13).}
\label{tab_rf_1_13}
\begin{tabularx}{\textwidth}{l X c c}
\toprule
\textbf{Código} & \textbf{Descrição do Requisito Funcional} & \textbf{Prioridade} & \textbf{Atores} \\
\midrule
\textbf{RF01} & O sistema deve restringir o cadastro e login de usuários estritamente ao domínio de e-mail institucional \texttt{@academico.ufs.br}, rejeitando qualquer outro domínio no backend e frontend. & Essencial & Todos \\
\midrule
\textbf{RF02} & O sistema deve gerar uma senha provisória aleatória no primeiro cadastro do usuário e forçar obrigatoriamente a troca de senha no primeiro login antes de liberar o acesso às funcionalidades. & Essencial & Todos \\
\midrule
\textbf{RF03} & O sistema deve permitir que qualquer usuário autenticado altere sua senha voluntariamente em seu perfil pessoal, mediante validação da senha atual. & Essencial & Todos \\
\midrule
\textbf{RF04} & O sistema deve gerenciar usuários contemplando estritamente os três perfis de acesso: Administrador (\texttt{ADMIN}), Professor (\texttt{PROFESSOR}) e Estudante (\texttt{STUDENT}). & Essencial & ADMIN \\
\midrule
\textbf{RF05} & O campo Matrícula deve ser de preenchimento obrigatório exclusivamente para discentes (\texttt{STUDENT}), permanecendo oculto e nulo para \texttt{ADMIN} e \texttt{PROFESSOR}. & Essencial & ADMIN \\
\midrule
\textbf{RF06} & O sistema deve prover CRUD completo de Módulos Assistenciais GTT com paginação de exatamente 10 registros por página, busca textual com debounce e validação de código único. & Essencial & ADMIN \\
\midrule
\textbf{RF07} & O sistema deve prover CRUD completo dos 53 Gatilhos Clínicos GTT, validando integridade referencial com o módulo assistencial e unicidade de código do gatilho. & Essencial & ADMIN \\
\midrule
\textbf{RF08} & O sistema deve fornecer um Catálogo Educacional de Gatilhos e Módulos para consulta livre de alunos e professores, com accordions contendo diretrizes e pistas clínicas. & Importante & Todos \\
\midrule
\textbf{RF09} & O sistema deve prover CRUD completo e Guia Clínico Interativo das Categorias de Gravidade de Dano NCC MERP (A a I), com distinção visual de dano real (E a I). & Essencial & Todos \\
\midrule
\textbf{RF10} & O sistema deve gerenciar Turmas Acadêmicas formatadas estritamente no padrão universal: \texttt{"Nome da Matéria - Turma - Período"} (ex: \texttt{Física 3 - T06 - 2026.2}). & Essencial & ADMIN \\
\midrule
\textbf{RF11} & Toda turma acadêmica deve possuir obrigatoriamente exatamente um Professor Titular vinculado, permitindo substituição atômica sem deixar a vaga docente em aberto. & Essencial & ADMIN \\
\midrule
\textbf{RF12} & O sistema deve permitir a matrícula e desvinculação de múltiplos estudantes na turma acadêmica, validando a existência e o perfil de discente do usuário. & Essencial & ADMIN \\
\midrule
\textbf{RF13} & O sistema deve bloquear o encerramento de uma turma acadêmica caso existam atividades com submissões pendentes de correção pelo professor. & Essencial & ADMIN \\
\bottomrule
\end{tabularx}
\fonte{Elaborado pelo autor (2026).}
\end{table}

\begin{table}[htb]
\centering
\small
\caption{Especificação dos Requisitos Funcionais (RF14 a RF25).}
\label{tab_rf_1_14}
\label{tab_rf_14_25}
\begin{tabularx}{\textwidth}{l X c c}
\toprule
\textbf{Código} & \textbf{Descrição do Requisito Funcional} & \textbf{Prioridade} & \textbf{Atores} \\
\midrule
\textbf{RF14} & O docente deve poder criar, editar e excluir Atividades Acadêmicas vinculadas à sua turma, definindo título, orientações didáticas e data/hora limite de entrega. & Essencial & PROFESSOR \\
\midrule
\textbf{RF15} & O sistema deve permitir a elaboração de Prontuários Fictícios Simulados em JSONB (dados do paciente, evolução multiprofissional, prescrições, exames e cirurgias). & Essencial & PROFESSOR \\
\midrule
\textbf{RF16} & A tela de resolução de atividade deve incorporar um cronômetro regressivo estrito de 20 minutos com controle de pausa/retomada e alertas visuais de proximidade de término. & Essencial & STUDENT \\
\midrule
\textbf{RF17} & O estudante deve selecionar os gatilhos identificados no caso clínico, indicando a presença de dano real, a categoria de gravidade NCC MERP (E a I) e a justificativa clínica. & Essencial & STUDENT \\
\midrule
\textbf{RF18} & O sistema deve prover interface interativa para construção do Diagrama de Causa e Efeito de Ishikawa, estruturado nos 6M em torno do problema central investigado. & Essencial & STUDENT \\
\midrule
\textbf{RF19} & O sistema deve prover interface para Matriz GUT com cálculo dinâmico em tempo real do Score de Priorização ($S_{GUT} = G \times U \times T$) para cada problema clínico relatado. & Essencial & STUDENT \\
\midrule
\textbf{RF20} & O sistema deve guiar o preenchimento da Matriz de Ação 5W2H com suporte a múltiplas linhas dinâmicas de contramedidas preventivas hospitalares. & Essencial & STUDENT \\
\midrule
\textbf{RF21} & O sistema deve prover campos estruturados para o Ciclo de Melhoria Contínua PDCA/PDSA (\textit{Plan}, \textit{Do}, \textit{Check}, \textit{Act}). & Essencial & STUDENT \\
\midrule
\textbf{RF22} & O sistema deve permitir o preenchimento da Matriz SWOT (Forças, Oportunidades, Fraquezas e Ameaças) e painel dinâmico de ideação por Brainstorming. & Essencial & STUDENT \\
\midrule
\textbf{RF23} & O sistema deve permitir a submissão individual da resolução da atividade pelo estudante antes do prazo final estipulado. & Essencial & STUDENT \\
\midrule
\textbf{RF24} & O docente deve visualizar o espelho completo da resolução do discente e registrar a avaliação atribuindo nota de 0,00 a 10,00 e parecer pedagógico formativo obrigatório. & Essencial & PROFESSOR \\
\midrule
\textbf{RF25} & O sistema deve consolidar o Dashboard Analítico da Turma com cálculo automático dos indicadores IHI ($TI_{1000}$, $TI_{100}$, $P_{EA}$) e métricas de desempenho pedagógico. & Essencial & PROFESSOR, ADMIN \\
\bottomrule
\end{tabularx}
\fonte{Elaborado pelo autor (2026).}
\end{table}

\section{Requisitos Não Funcionais (RNF01 a RNF10)}
\label{sec_requisitos_nao_funcionais}

Os Requisitos Não Funcionais estabelecem restrições arquiteturais, critérios de qualidade técnica, preceitos de usabilidade, padrões de segurança e desempenho que delimitam a solução (\autoref{tab_rnf}):

\begin{table}[htb]
\centering
\small
\caption{Especificação dos Requisitos Não Funcionais (RNF01 a RNF10).}
\label{tab_rnf}
\begin{tabularx}{\textwidth}{l X c}
\toprule
\textbf{Código} & \textbf{Descrição do Requisito Não Funcional} & \textbf{Dimensão} \\
\midrule
\textbf{RNF01} & \textbf{Resoluções de Tela Homologadas}: O layout deve ser projetado exclusivamente para ambientes desktop e estações de trabalho clínicas, com renderização perfeita e ergonomia em resoluções HD ($1280 \times 720$), HD+ ($1600 \times 900$), WUXGA ($1920 \times 1200$) e monitores Ultrawide 21:9 ($2560 \times 1080$ e $3440 \times 1440$). & Usabilidade e Design \\
\midrule
\textbf{RNF02} & \textbf{Desempenho da API}: O tempo médio de resposta dos endpoints REST no backend deve ser inferior a 200 milissegundos para 95\% das requisições sob condições normais de uso na rede acadêmica da UFS. & Desempenho \\
\midrule
\textbf{RNF03} & \textbf{Segurança e Autenticação}: A comunicação cliente-servidor deve utilizar autenticação \textit{stateless} com JSON Web Token (JWT) assinado com HMAC-SHA256, expiração controlada e senhas com hashing criptográfico adaptativo BCrypt. & Segurança \\
\midrule
\textbf{RNF04} & \textbf{Isolamento de Dados e LGPD}: O software não deve estabelecer conexões diretas com os bancos de dados de produção do hospital universitário, garantindo conformidade absoluta com a Lei Geral de Proteção de Dados (LGPD). & Privacidade e Bioética \\
\midrule
\textbf{RNF05} & \textbf{Padronização Universal de Paginação}: Todas as tabelas de listagem do sistema devem adotar paginação fixa de exatamente 10 registros por página, com controles padronizados de navegação e indicadores de contagem. & Usabilidade \\
\midrule
\textbf{RNF06} & \textbf{Ergonomia Visual e Paleta Cromática Hospitalar}: A interface deve adotar azuis clínicos profundos (\texttt{\#1E3A8A}, \texttt{\#2563EB}), verde esmeralda suave para confirmações, cinzas de baixo cansaço visual, âmbar para alertas e carmim cirúrgico para gravidades H e I. & Design Clínico \\
\midrule
\textbf{RNF07} & \textbf{Feedback e Indicador Global de Carregamento}: Todas as operações assíncronas devem exibir \textit{spinner} de carregamento sobreposto e notificações reativas \textit{Toast} para sucessos, alertas e erros. & Usabilidade \\
\midrule
\textbf{RNF08} & \textbf{Prevenção de Ações Destrutivas}: Ações irreversíveis (exclusão de turmas, atividades ou usuários e encerramento de turmas) devem obrigatoriamente exigir confirmação explícita através de diálogos modais dedicados. & Confiabilidade \\
\midrule
\textbf{RNF09} & \textbf{Critério Estrito de Testes Automatizados}: O projeto deve atingir e sustentar 100\% de cobertura de testes automatizados de ramos e linhas tanto no backend (JUnit 5 + Jacoco) quanto no frontend (Jest). & Confiabilidade \\
\midrule
\textbf{RNF10} & \textbf{Documentação Completa da Base de Código}: Todo o código deve ser documentado sem erros: Javadoc e SpringDoc OpenAPI 3 no backend; Compodoc no frontend Angular. & Manutenibilidade \\
\bottomrule
\end{tabularx}
\fonte{Elaborado pelo autor (2026).}
\end{table}
